%%%%%%%%%%%%%%%%%%%%%%%%%%%%%%%%%%%%%%%%%%%%%%%%%%%%%%%%%%%%%%%%%%%%%%%%%%%%%%%
%% Abstract --- tradução do resumo.
%%%%%%%%%%%%%%%%%%%%%%%%%%%%%%%%%%%%%%%%%%%%%%%%%%%%%%%%%%%%%%%%%%%%%%%%%%%%%%%

\begin{resumo}[Abstract]
\begin{otherlanguage*}{english}
%% Ponto decimal só neste ambiente; o resto do documento usa vírgula.
\sisetup{output-decimal-marker = {.}}
Vibration analysis lets maintenance decisions follow the measured condition of a motor, but the
cost of instrumentation limits its use to critical machines and leaves most motors in a plant
unmonitored. Low-cost sensor nodes could extend this coverage, yet the studies proposing them rarely
report the bandwidth and noise floor of the sensor. This work investigated which motor failure modes
a low-cost measurement chain, made up of a MEMS accelerometer sampled at \SI{1}{\kilo\hertz}, an ESP32
microcontroller and a LoRa link, can detect, and how much estimation decisions change the reported
performance. The aim was to measure detection by failure mode and to calibrate the threshold of an
Isolation Forest detector for a fixed false alarm rate. The experiments used \num{1951} test-rig
recordings from the public MAFAULDA dataset, decimated to \SI{1}{\kilo\hertz}, and synthesised signals,
with integration tested against an analytical value. On real data, at a \SI{2.0 \pm 1.3}{\percent}
false alarm rate, detection was \SI{75.0 \pm 3.5}{\percent} for imbalance, \SI{14.8}{\percent} and
\SI{20.7}{\percent} for bearing faults with no added mass and \SI{5.0}{\percent} and \SI{7.2}{\percent}
for misalignment; with the default threshold, the false alarm rate was \SI{31.2 \pm 4.2}{\percent}. On
synthesised signals, a cumulative sum without mean removal missed the RMS velocity by \SI{73}{\percent},
and by \SI{716}{\percent} with a bias within the sensor specification, whereas frequency-domain
integration came within \SI{0.01}{\percent} of the analytical value; the partial area under the ROC curve
fell by up to \SI{20.3}{\percent} between the \SI{20}{\kilo\hertz} and \SI{1}{\kilo\hertz} chains, against
\SI{3.7}{\percent} for the global area; and a supervised classifier reached \SI{99.0}{\percent} accuracy
with random window partitioning and \SI{63.2}{\percent} with partitioning by test run. \emandamento{In two
bench setups, with no motor and no induced fault, the device sampled at about \SI{999.8}{\hertz} over three
cycles and the receiver accepted \num{270} packets.} The results indicate that, with this detector and four
features, the chain can detect imbalance at a fixed false alarm rate and detects few bearing faults when no
imbalance is added. They also indicate that the performance of devices of this class can only be compared
when the bandwidth, threshold and validation protocol are reported.

\vspace{\onelineskip}
\noindent
\textbf{Keywords}: predictive maintenance; vibration analysis; anomaly detection;
embedded systems; low-cost instrumentation.
\end{otherlanguage*}
\end{resumo}
